\noindent\makebox[0pt][l]{364}\hfill{\small\textbf{BRIEFWECHSEL MIT BOHNENBERGER.}}\hfill\null\par\medskip

\begin{center}
[3].\\
\textsc{Gauss} an \textsc{Bohnenberger}.\quad Göttingen, 16.~November 1823.
\end{center}

\par\noindent\hspace*{\parindent}Mit Vergnügen habe ich aus Ihrem Briefe einiges von Ihren Messungen erfahren, und möchte nichts lieber, als dass diese Antwort Veranlassung geben möchte, etwas ausführlicheres darüber mitgetheilt zu erhalten.

\par\noindent\hspace*{\parindent}Was zuerst die Heliotrope betrifft, so habe ich zwei ganz verschiedene Arten anfertigen lassen, die auf ganz verschiedenen Principien beruhen; die zweite Einrichtung finde ich aber am vortheilhaftesten.\ .\ .\ .\ .\ .\quad Durch die Möglichkeit, wo es sonst das Terrain erlaubt, die grössten Dreiecke anzuwenden, überall gleich anfangen zu können, ohne erst die so viel Zeit und Geld kostenden Signale errichten zu müssen, wird die kleine Ausgabe [für die Heliotrope] vielfach erspart, obwohl dies der geringste Vortheil ist; die Messungen werden dadurch einer Schärfe fähig, auf die man bei Signalen und Kirchthürmen selten rechnen darf.
Meine schlechtesten Dreiecke (relativ gesprochen) sind die, worin Thürme die Zielpunkte waren.
So viel von den Heliotropen.

\par\noindent\hspace*{\parindent}Was die Messungen selbst betrifft, so wünsche ich nichts sehnlicher, als bald die Verbindung mit Ihnen zu haben.
Ich habe jetzt meine Dreiecke mit den kurhessischen zusammengehängt, wovon beiliegende Zeichnung Ihnen einen Begriff gibt~[\textsuperscript{*})].
Zwar ist von den hessischen Dreiecken nur erst ein Theil wirklich gemessen, aber schon genug, um alle Punkte bis zum Feldberg mit den meinigen wenigstens vorläufig zu verknüpfen.
Meine nördlicher liegenden Dreiecke kennen Sie aus \textsc{Schumachers} Astronomischen Nachrichten, I.~Nr.~24, die in diesem Jahr noch dazu gekommenen nördlichen Punkte werden Sie wenigstens für den Augenblick nicht interessiren.

\par\noindent\hspace*{\parindent}Ich vermuthe nun, dass hiedurch und die darmstädtschen und die bayerischen Dreiecke meine Messungen mit den Ihrigen verbunden sind, besitze aber von den letztern noch \textls[120]{gar keine}, von den erstern nur einige fragmentarische Angaben.
Sollten Sie in vollständigerm Besitze sein, so verpflichten Sie mich ausserordentlich durch Mittheilung, ebenso wie von Ihren eigenen Dreiecken.
Nur wünschte ich die Winkel zwar auf die Centra reducirt, aber ohne Fehlerausgleichung zu erhalten~\textsuperscript{**}).
Ich möchte gern alles nach gleichför-
\par\smallskip\noindent\rule{.24\textwidth}{.4pt}\par\smallskip
{\footnotesize
\noindent\hspace*{1.5em}[\textsuperscript{*}) Eine Zeichnung liegt der im Gauss-Archiv befindlichen Copie des Originals nicht bei.]\par
\noindent\hspace*{1.5em}\textsuperscript{**}) Übrigens allenfalls ohne alle weitern Rechnungsresultate.\par
}
\clearpage
\noindent\null\hfill{\small\textbf{ZUR HANNOVERSCHEN TRIANGULATION.}}\hfill\makebox[0pt][r]{365}\par\medskip
\noindent miger Methode berechnen.
Die Verknüpfung mit Tübingen, Mannheim und München interessirt mich um so dringender, da, \textls[120]{im Vertrauen gesagt}, die ersten Messungen \textsc{Schumachers} in Altona, wo er eine schöne kleine Sternwarte mit einem dem meinigen, \textsc{Soldner}schen und \textsc{Bessel}schen ganz gleichen \textsc{Reichenbach}schen Meridiankreise errichtet hat, eine Polhöhe gegeben haben, die \(5^{\prime\prime}\) von der von Göttingen durch die Dreiecke übertragenen differirt.
Bei der Rechnung habe ich die Dimensionen der Erde, die \textsc{Walbeck} aus dem Ensemble aller guten Gradmessungen abgeleitet hat, gebraucht: Abpl.\ \(\frac{1}{302{,}78}\).
Übrigens aber sind meine Rechnungsmethoden so gänzlich von den sonst angewandten verschieden, dass ich in einem Briefe Ihnen keinen Begriff davon geben kann.
Ich habe die Absicht, wenn meine Messungen erst vollendet sind, diese Methoden zu einem grössern Werke zu verarbeiten und durch Anwendung auf die hannoverschen und die damit zusammenhängenden Messungen zu erklären.
Unser Gouvernement ist geneigt, die hannoverschen Messungen weiter westlich auszudehnen, wodurch sie mit den \textsc{Krayenhoff}schen und dadurch mit den französischen und englischen in Zusammenhang kommen.
\textsc{Krayenhoffs} Messungen sind bekanntlich gedruckt; dies sollte mit allen guten Triangulirungen geschehen; die grossen Dreiecke gehören gewissermaassen der ganzen cultivirten Mit- und Nachwelt an, um so mehr, je mehr sie nach und nach unter sich in Zusammenhang kommen.
Die \textsc{Müffling}schen Dreiecke hängen zwar mit der französischen Gradmessung auch schon durch \textsc{Tranchots} Dreiecke zusammen; allein diese sind nicht gedruckt, also für das Publicum so gut wie gar nicht vorhanden, und General \textsc{von Müffling} selbst besitzt sie nur in ungenügender Form (vermuthlich die fatalen Chordenwinkel zu \(180^{0}\) schon abgeglichen).
Leider finde ich, dass die Menschen so wenig zur Communication geneigt sind; ich habe mir auf officielem diplomatischen und auf nicht officielem Privatwege viele Mühe gegeben, aus Paris die von \textsc{Epailly} 1804 und 1805 im Hannoverschen gemachten Messungen zu erhalten, aber nichts als Ausflüchte, eine blosse Namenangabe der Stationen und eine Zeichnung der Dreiecke erhalten, 2 oder 3 Zahlangaben nicht gerechnet, die, wie aus meinen Messungen folgt, entschieden grob unrichtig sind.
Lassen Sie uns eine Ausnahme davon machen.
Ich wiederhole nochmals meine Bitte um eine Mittheilung Ihrer Messungen, und um freundschaftliche Mittheilung dessen, was Sie von fremden mit den unsrigen zusammenhängenden haben, insofern ich es direct nicht erhalten kann, und erbiete mich gern ad reciproca.
\par\clearpage
\noindent\makebox[0pt][l]{366}\hfill{\small\textbf{BRIEFWECHSEL MIT BOHNENBERGER.}}\hfill\null\par\medskip
\noindent Wie schön wäre es, wenn einmal alle über Europa, von Schottland bis zum Banat und von Copenhagen bis Genua und Formentera, sich erstreckenden Messungen in Ein zusammenhängendes System gebracht werden könnten.
Ich möchte gern nach Kräften dazu vorbereiten, allein wenn man über die Mitte seines Lebens hinaus ist, muss man bei einem so ausgedehnten Gegenstand je eher je lieber anfangen.

\par\noindent\hspace*{\parindent}\textsc{Gerlings} Messungen sind mit einem 12-zölligen \textsc{Ertel}schen Theodolithen gemacht, ganz dem meinigen und \textsc{Schumacher}schen gleich\textsuperscript{*}).
Bei meinen Messungen habe ich gefunden, dass das, was ich in meiner Abhandlung in den neuesten Göttinger Commentationes »Theoria combinationis observationum etc.« den \textls[120]{mittlern} Fehler nenne, aus mehrern Stationen, gute und \textls[120]{weniger gute} Messungen durch einander gerechnet, etwa \(=\frac{3\mathord{\overset{\prime\prime}{,}}5}{\surd n}\) ist, \(n=\) Anzahl der Repetitionen.
Bei sehr fester Aufstellung, sehr günstiger (d.~i.\ nicht zitternder) Luft und ausschliesslich heliotropischen Zielpunkten ist er aber beträchtlich kleiner.
Meine sämmtlichen Messungen geben bisher 76 Hauptrichtungen (38 hin und 38 zurück) und aus der Ausgleichung der Fehler fand sich, dass der \textls[120]{mittlere} Fehler einer Hauptrichtung \(=0\mathord{\overset{\prime\prime}{,}}47\) war.
Es bilden sich daraus zusammen 26 Dreiecke, worin alle Winkel von mir gemessen sind; darunter mehrere, die gekreuzte Vierecke und Fünfecke geben, und die Ausgleichung ist ohne Willkür, ohne Auswählen und ohne Ausschliessen gemacht, nach strengen Gründen der Wahrscheinlichkeits-Rechnung, so dass zuletzt alles genau zu einander passt.
Solche gekreuzte Vierecke würden bei manchen Messungen ein trefflicher Probirstein sein, wo man findet, dass die Summen der Winkel zwar überall vortrefflich passen, so dass selten ein Dreieck viel über \(1^{\prime\prime}\) fehlt, wo aber jene Prüfung (die nicht in dem Grade à la portée von jedermann ist, wie das Berechnen eines sphärischen Excesses), wenn sie angewandt werden kann, zuweilen ganz entschieden zeigt, dass Fehler von \(2^{\prime\prime}\), \(3^{\prime\prime}\) oder darüber in einzelnen Winkeln vorhanden sind.
Der grösste Fehler in der Summe der noch nicht ausgeglichenen Winkel bei meinen 26 Dreiecken war \(2\mathord{\overset{\prime\prime}{,}}2\), wo eine Richtung auf den sehr schwer zu schneidenden und bei Sonnenschein nicht ganz phasenfreien Michaelisthurm in Hamburg ging; der nächst grösste \(1\mathord{\overset{\prime\prime}{,}}8\) in einem Dreiecke, wo auch eine Richtung auf einen äusserst schwer zu sehen-
\par\smallskip\noindent\rule{.24\textwidth}{.4pt}\par\smallskip
{\footnotesize\noindent\hspace*{1.5em}\textsuperscript{*}) Derselbe hat auch 3 Heliotrope, von Herrn \textsc{Rumpf}, nach der zweiten Einrichtung.\par}
\clearpage
\noindent\null\hfill{\small\textbf{ZUR HANNOVERSCHEN TRIANGULATION.}}\hfill\makebox[0pt][r]{367}\par\medskip
\noindent den nicht heliotropischen Zielpunkt ging.
Ich hatte gewünscht, die letztere betreffende Station noch einmal zu besuchen und die Richtung durch Heliotroplicht zu nehmen, konnte aber in diesem Jahre nicht mehr dazu kommen.
Das grosse Dreieck Hohehagen-Brocken-Inselsberg ist unter den 26 nicht begriffen, der Winkel auf dem Inselsberg ist von \textsc{Gerling} gemessen, und meinen Messungen auf dem Brocken war das Wetter sehr ungünstig, so dass ich den Inselsberg nur 15-mal habe schneiden können (1823); verbunden mit den 15 Schnitten von 1821, die weniger als \(1^{\prime\prime}\) von jenen differiren, geben sie aber doch einen vortrefflichen Schluss dieses grossen Dreiecks. --- In einem Lande wie Württemberg und Kurhessen, wo es so viele hohe Punkte gibt, ist das Messen ein Vergnügen, und die grössten Dreiecke leicht aufzufinden.
Ebenso im südlichen Theile des Hannoverschen.
Aber im nördlichen, der Lüneburger Heide, habe ich unsägliche Schwierigkeiten gehabt, und eine im vorigen Sommer nach Westen, gegen Bremen zu, unternommene Recognoscirung hat noch keine Resultate für die Möglichkeit nur leidlich guter Dreiecke gegeben.
Durch hohe Gerüste liessen sich die Schwierigkeiten zwar wohl überwinden, aber ich fürchte mich vor den grossen Kosten an Geld und Zeit und noch mehr vor der Einbusse solider Aufstellung.
Weiter südlich durch das Osnabrücksche nach Bentheim liesse sich vermuthlich leichter durchkommen, die Messungen würden dann aber dem grössten Theile nach über fremdes Gebiet gehen müssen.

\begin{center}\rule{.24\textwidth}{.4pt}\end{center}
